\documentclass[12pt,a4paper]{article}

\usepackage[utf8]{inputenc}
\usepackage[T1]{fontenc}
\usepackage[a4paper]{geometry}

\usepackage{pdflscape}
\usepackage{float}

\usepackage{amsthm}
\usepackage{amsfonts}
\usepackage{amsmath} %pakiet potrzebny do \eqref
\usepackage{amssymb} %pakiet potrzebny do \blacksquare i \mathbb{•}

\usepackage[polish]{babel}

\usepackage{graphicx}

\newtheorem{lemat}{Lemat}
\newtheorem{tw}{Twierdzenie}
\newtheorem{przyklad}{Przykład}
\theoremstyle{definition}
\newtheorem{df}{Definicja}

\usepackage[pdfauthor={Janisz},%
pdftitle={AM2},%
pagebackref=true,%
pdftex]{hyperref}
\hypersetup{colorlinks=false}


\pagestyle{plain}
\begin{document}
\title{ Analiza matematyczna 2}
\author{\vspace{-5ex}}
\date{\today}
\maketitle
\tableofcontents

\DeclareGraphicsExtensions{.pdf,.png,.jpg}
\begin{center}
\leavevmode

\includegraphics[width=1 in]{by-sa.png}
\end{center}
\label{fig:cc}
%insert a link to the licence and its description below
\scriptsize{Ten utwór jest dostępny na licencji  \href{http://creativecommons.org/licenses/by-sa/3.0/pl/}{Creative Commons Uznanie autorstwa-Na tych samych warunkach 3.0 Polska.}}


\pagebreak

\section{Całka Riemanna}
\begin{df}[Ciąg normalny podziałów]
Ciąg podziałów $(\pi(n))$ nazywamy normalnym, jeśli $\delta (\pi (n)) \xrightarrow{n\to\infty}0$
\end{df}
\begin{df}[Sumy całkowe]~\\
$\begin{array}{l l l l}
\text{Dolna suma całkowa:}& s_n = s_n(\pi_n) = \sum \limits_{i=1}^{k_n}m_i^{(n)}
(x_i^{(n)}-x_{i-1}^{(n)}) & m_i^{(n)} =
 \inf (f(x)) & x\in [x_{i-1}^{(n)} ,x_{i}^{(n)} ]\\
\text{Górna suma całkowa:}& S_n = S_n(\pi_n) = \sum_{i=1}^{k_n}M_i^{(n)}(x_i^{(n)}-x_{i-1}^{(n)}) & M_i^{(n)} = \sup (f(x)) & x\in [x_{i-1}^{(n)},x_{i}^{(n)}]\\
\text{Suma całkowa Riemanna:}& \sigma_n = \sigma_n(\pi_n) = \sum_{i=1}^{k_n}f(\xi_i^{(n)})(x_i^{(n)}-x_{i-1}^{(n)}) & & \xi\in [x_{i-1}^{(n)},x_{i}^{(n)}]
\end{array}$
\end{df}
\begin{df}[Funkcja całkowalna w sensie Riemanna]~\\
Funkcja f jest całkowalna w sensie Riemanna na $[a,b] \Leftrightarrow$ istnieje $\sigma\in\mathbb{R}$ taka,że dla dowolnego normalnego ciągu podziałów $(\pi_n)$ oraz dla dowolnego wartościowania tego ciągu $(\omega_n)$ mamy $\lim\limits_{n\to\infty}\sigma_n(\pi_n,\omega_0) = \lim\limits_{n\to\infty}\sum\limits_{i=1}^{k_n}f(\xi_i^{(n)})(x_i^{(n)}-x_{i-1}^{(n)})$, to $\sigma$ nazywamy całką Riemanna funkcji f na $[a,b]$ i oznaczamy $\int\limits_a^b(f(x)dx$
\end{df}
\begin{tw}
$f: [a,b] \rightarrow \mathbb{R}$ (ograniczona) jest całkowalna w sensie Riemanna na $[a,b] \Leftrightarrow s=S$
\end{tw}
\begin{tw}
Każda funkcja monotoniczna i ograniczona na $[a,b]$ jest całkowalna w sensie Riemanna na $[a,b]$
\end{tw}
\begin{tw}
Każda funkcja ciągła na $[a,b]$ jest całkowalna w sensie Riemanna na $[a,b]$
\end{tw}

\subsection{Własności całki Riemanna}

\begin{enumerate}
\item $f,g: [a,b] \to \mathbb{R}, \quad\{x\in[a,b]: f(x) \neq g(x)\}$ jest zbiorem skończonym $\Rightarrow$ f jest całkowalna na $[a,b]\Leftrightarrow$ g jest całkowalna na $[a,b]$. W przypadku całkowalności $\int\limits_a^bf(x)dx=\int\limits_a^bg(x)dx$
\item $ \int\limits_{a}^{b} (\alpha f(x) + \beta g(x))dx = \alpha \int\limits_a^b f(x)dx + \beta \int\limits_a^bg(x)dx $
\item $f: [a,b]\to\mathbb{R}$ całkowalna na $[a,b] \Rightarrow$ f jest całkowalna na każdym pod przedziale $[a,b]$
\item $f: [a,b]\to\mathbb{R}, ~ c\in [a,b]$, f jest całkowalna na $[a,c]$ oraz $[c,b] \Rightarrow$ f jest całkowalna na $[a,b]$ oraz $\int\limits_a^bf(x)dx = \int\limits_a^cf(x)dx + \int\limits_c^bf(x)dx$
\item $f: [a,b] \to \mathbb{R}$ całkowalna na $[a,b]$, $\forall x\in[a,b] ~~ f(x) \geq 0 \Rightarrow \int\limits_a^bf(x)dx \geq 0$
\item $f,g: [a,b] \to \mathbb{R}$ całkowalne na $[a,b]$, $\forall x\in[a,b] ~~ f(x) \leqslant g(x) \Rightarrow \int\limits_a^bf(x)dx \leqslant \int\limits_a^bg(x)dx$
\item $f: [a,b] \to \mathbb{R}$ całkowalna na $[a,b] \Rightarrow$ $|f|$ też jest całkowalna na $[a,b]$ oraz $|\int\limits_a^bf(x)dx| \leqslant \int\limits_a^b|f(x)|dx$
\item $f: [a,b] \to \mathbb{R}$ całkowalna na $[a,b] \Rightarrow \int\limits_a^bf(x)dx \leqslant \sup\limits_{x\in[a,b]}f(x)\cdot(b-a)$
\item Podstawowy wzór rachunku całkowego\\
$f: [a,b] \to \mathbb{R}$ funkcja ciągła, $\phi(x)$ --- dowolna funkcja pierwotna dla $f(x)$, to znaczy $\phi'(x)=f(x) \Rightarrow \int\limits_a^bf(x)dx = \phi(b)-\phi(a)$
\item Całkowanie przez podstawienie\\
$
f: [a,b]\to\mathbb{R}$, funkcja ciągła, $g: [a,b]\to\mathbb{R},~ g\in C^1([a,b])\\
\alpha = g(a),~ \beta =g(b)\\
\int\limits_a^bf(g(x))g'(x)dx = \int\limits_\alpha^\beta f(t)dt, \text{ gdzie }t = g(x)
$
\begin{proof}~\\
	Niech $\phi(x)$ będzie funkcją pierwotną dla $f(x)$, to znaczy
	 $\phi'(x) = f(x), ~ P=\int\limits_\alpha^\beta f(t)dt = \phi(\beta)-\phi(\alpha)$\\
$[\phi(g(x))]' = \phi'(g(x))g'(x) = f(g(x))g'(x) \Rightarrow \phi(g(x)) \text{ to funkcja pierwotna funkcji } f(g(x))g'(x)$\\
$L = \int\limits_a^bf(g(x))g'(x)dx = \phi(g(b))-\phi(g(a)) = \phi(\beta)-\phi(\alpha)\\
L = P$
\end{proof}
\item Całkowanie przez części\\
$
u,v: [a,b]\to\mathbb{R} \text{ --- funkcje klasy } C^1\\
\int\limits_a^bu(x)v'(x)dx = [u(x)v(x)]_a^b - \int\limits_a^bu'(x)v(x)dx
$
\begin{proof}~\\
$
	[u(x)v(x)]' = u'(x)v(x)+u(x)v'(x)\\
	\int\limits_a^b[u(x)v(x)]'dx = \int\limits_a^b(u'(x)v(x)+u(x)v'(x))dx = 
	\int\limits_a^b u'(x)v(x)dx + \int\limits_a^bu(x)v'(x)dx \quad(*)\\
	u(x)v(x) \text{ to funkcja pierwotna } [u(x)v(x)]'\\
	\int\limits_a^b[u(x)v(x)]'dx u(b)v(b) - u(a)v(a) = [u(x)v(x)]_a^b \quad(**)\\
	(*):(**) \Rightarrow [u(x)v(x)]_a^b = \int\limits_a^bu'(x)v(x)dx + \int\limits_a^b u(x)v'(x)dx\\
	\int\limits_a^bu'(x)v(x)dx = [u(x)v(x)]_a^b-\int\limits_a^bu'(x)v(x)dx	
$
\end{proof}
\item Twierdzenie o wartości średniej rachunku całkowego
$
f,g: [a,b]\to\mathbb{R} \text{ --- funkcje ciągłe, g jest nieujemna (niedodatnia)}\\
 \exists \xi\in [a,b] \int\limits_a^bf(x)g(x)dx = f(\xi)\int\limits_a^bg(x)dx
$
\end{enumerate}

\subsection{Całki niewłaściwe}

\begin{df}[Całka niewłaściwa I rodzaju]~\\
$f: [a, \infty) \to \mathbb{R} \text{ całkowalna na } [\alpha, \beta] ~ \forall \beta > \alpha$ oraz istnieje granica $\lim\limits_{\beta\to\infty} \int\limits_\alpha^\beta f(x)dx \Rightarrow$ granicę tę nazywamy całka niewłaściwą pierwszego rodzaju i oznaczamy $\int\limits_\alpha^\infty f(x)dx := \lim\limits_{\beta\to\infty} \int\limits_\alpha^\beta f(x)dx$\\
Ponadto jeśli granica ta istnieje i jest skończona to całkę niewłaściwą nazywamy zbieżną. Natomiast w pozostałych przypadkach całkę niewłaściwą nazywamy rozbieżną.
\end{df}

\begin{df}[Całka niewłaściwa II rodzaju]~\\
$f: [a, b) \to \mathbb{R} \text{ całkowalna na } [\alpha, \beta] \quad \forall a < \beta < b$ oraz istnieje granica $\lim\limits_{\beta\to b^-} \int\limits_a^\beta f(x)dx \Rightarrow$ granicę tę nazywamy całka niewłaściwą drugiego rodzaju i oznaczamy $\int\limits_a^b f(x)dx := \lim\limits_{\beta\to b^-} \int\limits_\alpha^\beta f(x)dx$\\
Ponadto jeśli granica ta istnieje i jest skończona to całkę niewłaściwą nazywamy zbieżną. Natomiast w pozostałych przypadkach całkę niewłaściwą nazywamy rozbieżną.
\end{df}

\subsection{Kryteria zbieżności}

\begin{tw}[Kryterium porównawcze]~\\
$f,g: [a,b) \to \mathbb{R} \text{ całkowalna na } [\alpha, \beta] \quad \forall a < \beta < b$ oraz
$\forall x\in [a,b) \quad 0 \leqslant f(x) \leqslant g(x)$ Wtedy:
\begin{itemize}
\item $\int\limits_a^bg(x)dx$ jest zbieżna $\Rightarrow \int\limits_a^bg(x)dx$ też jest zbieżna
\item $\int\limits_a^bf(x)dx$ jest rozbieżna $\Rightarrow \int\limits_a^bg(x)dx$ też jest rozbieżna
\end{itemize}
\end{tw}

\begin{tw}
$f: [a,b) \to \mathbb{R} \text{ całkowalna na } [\alpha, \beta] \quad \forall a < \beta < b$ oraz
$\int\limits_a^b |f(x)|dx$ jest zbieżna, to $\int\limits_a^b f(x)dx$ też jest zbieżna. Mówimy wtedy, że jest zbieżna bezwzględnie.
\end{tw} 

\subsection{Zastosowania geometryczne całki Riemanna}

\begin{tw}[Pole zbioru płaskiego]~\\
$f: [a,b] \to \mathbb{R}$ funkcja nieujemna i ciągła, $D = \{(x,y) \in \mathbb{R}^2: x\in [a,b[, y\in [0, f(x)]$
pole $D = |D| = \int\limits_a^b f(x)dx$
\end{tw}

\begin{tw}[Długość łuku]~\\
Niech $f: [a,b] \to \mathbb{R}$ -- funkcja klasy $C^1$. Wówczas długość łuku opisanego równaniem $y=f(x), x\in [a,b]$ dana jest wzorem $L = \int\limits_a^b \sqrt{1 + \left[f'(x)\right]^2}dx$
\end{tw}

\begin{tw}[Objętość bryły]~\\
Niech $f: [a,b] \to \mathbb{R}$ -- funkcja klasy $C^1$ oraz $V$ oznacza bryłę powstałą poprzez obrót
 dookoła osi OX krzywej $y = f(x), x\in [a,b]$. Wówczas objętość $V$ dana jest wzorem $V = \pi\int\limits_a^b f^2(x)dx$
\end{tw}


\section{Szeregi liczbowe}

\begin{df}[sumy szeregu liczbowego]~\\
Jeśli istnieje skończona lub nie $\lim\limits_{n\to\infty}S_n$, to nazywamy ją sumą szeregu $\sum\limits_{n=1}^\infty a_n$ i zapisujemy $\sum\limits_{n=1}^\infty a_n=\lim\limits_{n\to\infty}S_n$.
Jeśli $\lim\limits_{n\to\infty}S_n$ jest skończona, to szereg $\sum\limits_{n=1}^\infty a_n$ nazywamy zbieżnym; w pozostałych przypadkach (to znaczy gdy granica jest nieskończona lub nie istnieje) szereg ten nazywamy rozbieżnym
\end{df}

\begin{tw}[Warunek konieczny zbieżności szeregu]~\\
$\sum\limits_{n=1}^\infty a_n$ jest zbieżny $\Rightarrow \lim\limits_{n\to\infty}a_n = 0$
\begin{proof}
Zakładamy, że $\sum\limits_{n=1}^\infty a_n$ jest zbieżny $\Rightarrow (S_n)$ jest zbieżny; oznaczamy 
$S_n \xrightarrow{n\to\infty} S$\\
$a_n = S_n - S_{n-1} \xrightarrow{n\to\infty} S - S = 0$ 
\end{proof}
\begin{przyklad}[Powyższy warunek nie jest warunkiem dostatecznym]~\\
$\sum\limits_{n=1}^\infty \frac{1}{n}$ nie jest zbieżny, mimo że $\lim\limits_{n\to\infty} \frac{1}{n} = 0$
\end{przyklad}
\end{tw}

\begin{tw}[Warunek Cauch'ego zbieżności szeregu]~\\
Szereg $\sum\limits_{n=1}^\infty a_n$ jest zbieżny $\Leftrightarrow$ spełnia warunek Cauch'ego, to znaczy
$$
\forall \varepsilon > 0 ~~\exists N ~~\forall m > n > N \quad |a_{n+1} + a_{n+2} + \dots + a_m| < \varepsilon
$$
\end{tw}

\begin{tw}[O mnożeniu szeregu przez stałą]~\\
$\sum\limits_{n=1}^\infty a_n$ jest zbieżny i $\lambda\in\mathbb{R}$ wówczas szereg $\sum\limits_{n=1}^\infty \lambda a_n$ jest zbieżny i $\sum\limits_{n=1}^\infty \lambda a_n = \lambda \sum\limits_{n=1}^\infty a_n$
\end{tw}
\begin{tw}[O dodawaniu i odejmowaniu szeregów]~\\
$\sum\limits_{n=1}^\infty a_n, \sum\limits_{n=1}^\infty b_n$ są zbieżne wówczas szereg $\sum\limits_{n=1}^\infty a_n \pm \sum\limits_{n=1}^\infty b_n$ jest zbieżny i $\sum\limits_{n=1}^\infty a_n \pm \sum\limits_{n=1}^\infty b_n = \sum\limits_{n=1}^\infty (a_n \pm b_n)$
\end{tw}

\subsection{Kryteria zbieżności szeregów}

\begin{tw}[Kryterium porównawcze]~\\
$\forall n\in\mathbb{N} \quad 0 \leqslant a_n \leqslant b_n$, wówczas
\begin{itemize}
	\item $\sum\limits_{n=1}^\infty b_n$ jest zbieżny $\Rightarrow \sum\limits_{n=1}^\infty a_n$ jest zbieżny
	\item $\sum\limits_{n=1}^\infty a_n$ jest rozbieżny $\Rightarrow \sum\limits_{n=1}^\infty b_n$ jest rozbieżny
\end{itemize}
\end{tw}

\begin{tw}[Kryterium d'Alemberta]~\\
Niech $\forall n\in\mathbb{N} \quad a_n > 0$ i istnieje granica 
$\lim\limits_{n\to\infty} \frac{a_{n+1}}{a_n} = g$, wówczas
\begin{itemize}
	\item $g < 1 \Rightarrow \sum\limits_{n=1}^\infty a_n < \infty$
	\item $g > 1 \Rightarrow \sum\limits_{n=1}^\infty a_n = \infty$
	\item $g = 1 \Rightarrow $?
\end{itemize}
\end{tw}

\begin{tw}[Kryterium Cauch'ego]~\\
Niech $\forall n\in\mathbb{N} \quad a_n \geqslant 0$ i oznaczamy $g = \limsup\limits_{n\to\infty} \sqrt[n]{a_n}$, wówczas
\begin{itemize}
	\item $g < 1 \Rightarrow \sum\limits_{n=1}^\infty a_n < \infty$
	\item $g > 1 \Rightarrow \sum\limits_{n=1}^\infty a_n = \infty$
	\item $g = 1 \Rightarrow $?
\end{itemize}
\end{tw}

\begin{tw}[Kryterium całkowe zbieżności szeregu]~\\
$f: [1, \infty) \to \mathbb{R}$ funkcja nieujemna i nierosnąca. Wtedy 
$\sum\limits_{n=1}^\infty f(n)$ jest zbieżny $\Leftrightarrow \int\limits_{n=1}^\infty f(x)dx$ jest zbieżna.
\end{tw}

\begin{tw}[Kryterium Dirichleta]~\\

$\left.
\begin{array}{l}
	(a_n) \text{ to ciąg nierosnący i taki że } \lim\limits_{n\to\infty} a_n = 0\\
	(b_n) \text{ to ciąg taki że ciąg sum częściowych jest ograniczony }\footnotemark[1]
\end{array} 
\right\rbrace \Rightarrow \sum\limits_{n=1}^\infty a_nb_n$ jest zbieżny
\footnotetext[1]{ to znaczy $\exists M\in\mathbb{R} ~~ \forall n\in\mathbb{N} \quad |b_1 + b_2 + \dots b_n| \leqslant M$}
\end{tw}

\begin{tw}[Kryterium Leibniza]~\\
$(a_n)$ to ciąg nierosnący i taki, że $\lim\limits_{n\to\infty} a_n = 0 \Rightarrow \sum\limits_{n=1}^\infty a_n(-1)^{n+1} = a-1 - a_2 + a_3 - a_4 + \dots$ jest zbieżny
\begin{proof}
Niech $b_n = (-1)^{n+1}$, wówczas $b_1 + b_2 + \dots + b_n = 
	\begin{cases}
		1 \quad n=2k+1\\
		0 \quad n=2k\\
 	\end{cases}
\Rightarrow \forall n\in\mathbb{N} ~~ 0\leqslant b_1 + b_2 + \dots + b_n \leqslant 1$ to znaczy ciąg $(b_1 + b_2 + \dots + b_n)$ jest ograniczony. Są spełnione założenia kryterium Dirichleta $\Rightarrow \sum\limits_{n=1}^\infty a_nb_n = \sum\limits_{n=1}^\infty a_n(-1)^{n+1}$ jest zbieżny. 
\end{proof}
\end{tw}

\begin{df}
Szereg $\sum\limits_{n=1}^\infty a_n$ jest zbieżny bezwzględnie $\Leftrightarrow \sum\limits_{n=1}^\infty |a_n|$ jest zbieżny.
\end{df}

\begin{df}
Szereg który jest zbieżny ale nie jest zbieżny bezwzględnie nazywamy zbieżnym warunkowo.
\end{df}

\begin{tw}
Szereg $\sum\limits_{n=1}^\infty a_n$ jest zbieżny bezwzględnie $\Rightarrow$ Szereg $\sum\limits_{n=1}^\infty a_n$ jest zbieżny i $|\sum\limits_{n=1}^\infty a_n| \leqslant \sum\limits_{n=1}^\infty |a_n|$ 
\end{tw}

%dodać przykłady i wyjaśnieniea




\section{Ciągi i szeregi funkcyjne}
\subsection{Ciągi funkcyjne}
\begin{df}[Punktowej zbieżności]~\\
Mówimy, że ciąg funkcyjny $(f_n)$ jest zbieżny do funkcji $f$ (punktowo)
 $$\Leftrightarrow \forall x\in A \quad \lim\limits_{n\to\infty}\underbrace{f_n(x)}_{\text{ciąg liczbowy}} = \underbrace{f(x)}_{\text{liczba}} \Leftrightarrow \forall x\in A ~ \forall\varepsilon > 0 ~~ \exists N = N(\varepsilon, x) ~ \forall n \geqslant N \quad |f_n(x)-f(x)| < \varepsilon$$, oznaczamy $f_n \rightarrow f$
\end{df}

\begin{df}[Jednostajnej zbieżności]~\\
Mówimy, że ciąg funkcyjny $(f_n)$ jest zbieżny do funkcji $f$ (jednostajnie) $$\Leftrightarrow \forall\varepsilon > 0 ~~ \exists N = N(\varepsilon, x) ~ \forall n \geqslant N ~ \forall x\in A ~ \quad |f_n(x)-f(x)| < \varepsilon$$, oznaczamy $f_n \rightrightarrows f$
\end{df}

\begin{tw}[Warunek równoważny zbieżności jednostajnej ciągu funkcyjnego]~\\
$f, f_n: A \to \mathbb{R}$. Wówczas $f_n \rightrightarrows f \Leftrightarrow \lim\limits_{n\to\infty} \sup\limits_{x\in A} |f_n(x)-f(x)| = 0$
\end{tw}

\begin{przyklad}[ciągu który jest zbieżny punktowo ale nie jednostajnie]
$f_n: [0,1] \to \mathbb{R}, \quad f_n(x) = x^n$ zbieżność punktowa
$
	\lim\limits_{n\to\infty}f_n(x) = \lim\limits_{n\to\infty}x^n = 
	\begin{cases}
	0 \quad x\in[0,1)\\
	1 \quad x=1
	\end{cases}
$
\end{przyklad}

\begin{df}[Warunek Cauch'ego dla zbieżności punktowej i jednostajnej]~\\
$\forall x\in A f_n \rightarrow f \Leftrightarrow \forall x\in A$ ciąg $(f_n(x))$ jest zbieżny $\Leftrightarrow \forall x\in A \text{ ciąg } (f_n(x))$ spełnia warunek Cauch'ego, to znaczy 
$\forall x\in A ~~ \forall \varepsilon>0 ~~ \exists N ~ \forall n,m \geqslant N \quad |f_n(x) - f_m(x)| < \varepsilon$ 
\end{df}

\begin{tw}
Ciąg funkcyjny $(f_n)$ jest zbieżny jednostajnie na A $\Leftrightarrow$ spełnia warunek Cauch'ego zbieżności jednostajnej.
\end{tw}

\begin{tw}[o ciągłości granicy ciągu funkcyjnego]~\\
$f, f_n: A\to \mathbb{R}, f_n \rightrightarrows_A f$, funkcje $f_n$ są ciągłe w punkcie $a\in A ~~ \forall n\in N \Rightarrow f$ jest ciągła w punkcie $a$
\end{tw}

\begin{tw}[o różniczkowaniu granicy ciągu funkcyjnego]~\\
$A \subset \mathbb{R} \text{ -- przedział }$ są różniczkowalne w każdym punkcie przedziału A. Ciąg $(f_n')$ jest z zbieżny jednostajnie na A, czyli $\exists x_0 \in A  \quad (f_n'(x_0))$ jest zbieżny $\Rightarrow$
\begin{itemize}
	\item ciąg $(f_n)$ jest jednostajnie zbieżny na A do pewnej funkcji granicznej
	\item funkcja graniczna $f$ jest różniczkowalna na A i $\forall x\in A ~~ f'(x) = \lim\limits_{n\to\infty} f_n'(x) = (\lim\limits_{n\to\infty} f_n(x))'$
\end{itemize}
\end{tw}

\begin{tw}[o całkowaniu granicy ciągu funkcyjnego]~\\
$A \text{ -- przedział } \subset \mathbb{R}, f, f_n \in C(A), f_n \rightrightarrows_A f \Rightarrow \forall a,b \in A \quad \int\limits_a^bf(x)dx = \lim\limits_{n\to\infty}\int\limits_a^bf_n(x)dx$ 
\end{tw}

\subsection{Szeregi funkcyjne}
% dodać definicje
\begin{tw}[Warunek Cauch'ego zbieżności jednostajnej]~\\
Szereg $\sum\limits_{n=1}^\infty a_n$ jest zbieżny jednostajnie $\Leftrightarrow$ spełnia warunek Cauchego jednostajnej zbieżności szeregu funkcyjnego.
$\forall \varepsilon ~~\exists N ~ \forall m>n>N ~~ \forall x\in A \quad |a_{n+1}(x) + \dots + u_m(x)| < \varepsilon $
\end{tw}

\begin{tw}[Kryterium Weierstrassa]~\\
Jeśli istnieje ciąg liczbowy $(a_n)$ taki, że $\forall n\in \mathbb{N} ~~ \forall x\in A \quad |u_n(x)|\leqslant a_n$ i $\sum\limits_{n=1}^\infty a_n$ jest zbieżny, to szereg funkcyjny $\sum\limits_{n=1}^\infty u_n$ jest zbieżny jednostajnie ( i bezwzględnie)
\begin{przyklad}
Szereg $\sum\limits_{n=1}^\infty \frac{sin(nx)}{2^n}$ jest zbieżny jednostajnie na $\mathbb{R}$ bo
$\forall n\in \mathbb{N} ~~ \forall x\in\mathbb{R} \quad \frac{sin(nx)}{2^n} \leqslant \frac{1}{2^n} $ i $\sum\limits_{n=1}^\infty \frac{1}{2^n}$ jest zbieżny jako ciąg geometryczny.
\end{przyklad}
\end{tw}

\begin{tw}[o ciągłości sumy szeregu funkcyjnego]~\\
$S, u_n: A \to \mathbb{R}, ~~ \sum\limits_{n=1}^\infty u_n$ jest zbieżny jednostajnie do funkcji S, funkcje $u_n$ są ciągłe na $A~ \forall n\in\mathbb{N} \Rightarrow S=\sum\limits_{n=1}^\infty u_n$ też jest ciągła na $A$.
\end{tw}

\begin{tw}[o różniczkowaniu sumy szeregu funkcyjnego]~\\
$A$ -- przedział $\subset \mathbb{R}$\\
$u_n: A \to \mathbb{R}$ są różniczkowalne w każdym punkcie przedziału $A$\\
$\sum\limits_{n=1}^\infty u_n'$ jest zbieżny jednostajnie na $A$\\
$\exists x_0\in A ~~ \sum\limits_{n=1}^\infty u_n(x_0)$ jest zbieżny $\Rightarrow$
\begin{itemize}
	\item $\sum\limits_{n=1}^\infty u_n$ jest zbieżny jednostajnie
	\item $\sum\limits_{n=1}^\infty u_n$ jest różniczkowalna na $A$ i $(\sum\limits_{n=1}^\infty u_n)' = \sum\limits_{n=1}^\infty u_n'$
\end{itemize}
\end{tw}

\begin{tw}[o całkowaniu szeregu funkcyjnego]~\\
$A$ -- przedział $\subset \mathbb{R}$\\
$u_n: A \to \mathbb{R}$ są ciągłe na $A$\\
$\sum\limits_{n=1}^\infty u_n$ jest zbieżny jednostajnie
$ \Rightarrow \forall a,b\in A ~~ \int\limits_a^b (\sum\limits_{n=1}^\infty u_n(x))dx = \lim\limits_{n\to\infty}\int\limits_a^b (u_1(x) + u_2(x) + \dots + u_n(x))dx =  \sum\limits_{n=1}^\infty \int\limits_a^b u_n(x)dx$
\end{tw}

\subsection{Szeregi potęgowe}
\begin{df}
Szeregiem potęgowym nazywamy szereg funkcyjny postaci $\sum\limits_{n=0}^\infty a_nx^n$, gdzie $a_n\in\mathbb{R}$ dla $n\in\mathbb{N}$ 
\end{df}

\begin{tw}[D'Alemberta]~\\
Jeśli istnieje granica $\lim\limits_{n\to\infty} |\frac{a_{n+1}}{a_n}| = \lambda $, to promień zbieżności szeregu potęgowego $\sum\limits_{n=0}^\infty a_nx^n$ dany jest wzorem:
$ R = 
	\begin{cases}
		\frac{1}{\lambda}, & \lambda \in (0,\infty)\\
		+\infty, & \lambda = 0\\
		0, & \lambda = \infty
	\end{cases} $
\end{tw}

\begin{tw}[Cauch'ego-Hadamarda]~\\
Promień zbieżności szeregu potęgowego $\sum\limits_{n=0}^\infty a_nx^n$ dany jest wzorem 
$ R = 
	\begin{cases}
		\frac{1}{\lambda}, & \lambda \in (0,\infty)\\
		+\infty, & \lambda = 0\\
		0, & \lambda = \infty
	\end{cases} $, gdzie $\lambda = \limsup\limits_{n\to\infty}\sqrt[n]{|a_n|}$
\end{tw}

\begin{przyklad}[Szeregu o promieniu zbieżności $= 7$]
$\sum\limits_{n=1}^\infty \frac{1}{n7^n}x^n$
\end{przyklad}
\begin{przyklad}[Szeregu zbieżnego tylko dla $x=0$]
$\sum\limits_{n=1}^\infty n^nx^n$
\end{przyklad}
\begin{przyklad}[Szeregu zbieżnego tylko dla $x=5$]
$\sum\limits_{n=1}^\infty n^n(5-x)^n$
\end{przyklad}
\begin{przyklad}[Szeregu zbieżnego $\forall x\in\mathbb{R}$]
$\sum\limits_{n=1}^\infty \frac{x^n}{n!}$
\end{przyklad}

\begin{tw}[o ciągłości sumy szeregu potęgowego]~\\
Niech promień zbieżności $R$ szeregu $\sum\limits_{n=0}^\infty a_nx^n$ będzie dodatni. Wówczas funkcja $f(x)=\sum\limits_{n=0}^\infty a_nx^n$ jest ciągła na $(-R, R)$
\end{tw}

\begin{tw}[Abela]
Szereg potęgowy jest funkcją ciągłą w każdym punkcie, w którym jest zbieżny (w punktach końcowych przedziału mówimy o ciągłości jednostronnej) 
\end{tw}

\begin{tw}[o różniczkowaniu szeregu potęgowego]~\\
\begin{itemize}
	\item Oba szeregi $\sum\limits_{n=0}^\infty a_nx^n$ i $\sum\limits_{n=0}^\infty (a_nx^n)'$ mają te same promienie zbieżności
	\item Jeśli $R > 0$, to $f(x) = \sum\limits_{n=0}^\infty a_nx^n$ jest różniczkowalna na $(-R, R)$ i $f'(x) = (\sum\limits_{n=0}^\infty a_nx^n)' = \sum\limits_{n=0}^\infty (a_nx^n)' \quad \forall x\in (-R,R)$
\end{itemize}
\end{tw}

\begin{tw}[o całkowaniu szeregu potęgowego]~\\
\begin{itemize}
	\item Oba szeregi $\sum\limits_{n=0}^\infty a_nx^n$ i $\sum\limits_{n=0}^\infty \int\limits_0^x (a_nt^n)dt$ mają te same promienie zbieżności
	\item Jeśli $R > 0$, to $f(t) = \sum\limits_{n=0}^\infty a_nt^n$ jest całkowalna na $(0, x)$ i $\int\limits_0^x f(t)dt = \int\limits_0^x(\sum\limits_{n=0}^\infty a_nt^n)dt \quad \forall x\in (-R,R)$
\end{itemize}
\end{tw}

\subsection{Szereg Taylora i Maclaurina}
\begin{df}[szeregu Taylora]~\\
Niech $f\in C^\infty ((x_0-\delta, x_0+\delta))$ wtedy szereg potęgowy $\sum\limits_{n=0}^\infty \frac{f^{(n)}(x_0)}{n!}(x-x_0)^n$ nazywamy szeregiem Taylora o środku w punkcie $x_0$ dla funkcji $f$.
\end{df}

\begin{df}[szeregu Maclaurina]~\\
Niech $f\in C^\infty ((-\delta, \delta))$ wtedy szereg potęgowy $\sum\limits_{n=0}^\infty \frac{f^{(n)}(0)}{n!}x^n$ nazywamy szeregiem Maclaurina dla funkcji $f$.
\end{df}

\begin{przyklad}[Rozwinięć niektórych funkcji w szereg Macalurina]~\\
\begin{itemize}
	\item $\frac{1}{1-x} = \sum\limits_{n=0}^\infty x^n \quad |x| < 1$
	\item $e^x = \sum\limits_{n=0}^\infty \frac{x^n}{n!}$
	\item $sinx = \sum\limits_{n=0}^\infty \frac{(-1)^n}{(2n+1)!}x^{2n+1}$
	\item $cosx = \sum\limits_{n=0}^\infty \frac{(-1)^n}{(2n)!}x^{2n}$
\end{itemize}
\end{przyklad}

\section{Przestrzenie metryczne i unormowane}
\begin{df}
Przestrzenią metryczną nazywamy parę $(X, \rho)$, gdzie $X$ to niepusty zbiór, a $\rho$ to metryka w tym zbiorze. Elementy $X$ nazywamy punktami zaś $\rho (x,y)$ odległością między $x$ i $y$ 
\end{df}
\begin{przyklad}[Metryka naturalna (euklidesowa)]
$X = \mathbb{R}^2 ~\rho((x_1,x_2),(y_1, y_2)) = \sqrt{(x_1-y_1)^2+(x_2-y_2)^2}$
\end{przyklad}
\begin{przyklad}[Metryka dyskretna]
$X \text{-- dowolny zbiór niepusty} \quad \rho(x,y) = 
\begin{cases}
0 \text{ dla } x = y\\
1 \text{ dla } x \neq y
\end{cases}$
\end{przyklad}
\begin{przyklad}[Metryka taksówkowa (miejska)]
$X = \mathbb{R}^n ~\rho(x, y) = \sum\limits_{k=1}^n |x_k-y_k|$
\end{przyklad}

\begin{df}
Kulą (otwartą) o środku w punkcie $x_0$ i promieniu $r$ w przestrzeni metrycznej $(X, \rho)$ nazywamy zbiór $$K(x_0, r) = \{x\in X: \rho(x,x_0) < r\}$$
\end{df}
\begin{df}
Kulą domkniętą o środku w punkcie $x_0$ i promieniu $r$ w przestrzeni metrycznej $(X, \rho)$ nazywamy zbiór $$\overline{K}(x_0, r) = \{x\in X: \rho(x,x_0) \leqslant r\}$$
\end{df}
\begin{df}
Sferą o środku w punkcie $x_0$ i promieniu $r$ w przestrzeni metrycznej $(X, \rho)$ nazywamy zbiór $$S(x_0, r) = \{x\in X: \rho(x,x_0) = r\}$$
\end{df}

\begin{przyklad}[Kula w metryce dyskretnej]
$$K(x_0, r) = \{x\in X: \rho(x,x_0) < r\} = \begin{cases}
	\{x_0\} &\text{ gdy } r\in (0,1]\\
	X &\text{ gdy } r\in (1, \infty)
\end{cases}$$
\end{przyklad}
\begin{przyklad}[Kula domknięta w metryce dyskretnej]
$$\overline{K}(x_0, r) = \{x\in X: \rho(x,x_0) \leqslant r\} = \begin{cases}
	\{x_0\} &\text{ gdy } r\in (0,1)\\
	X &\text{ gdy } r\in [1, \infty)
\end{cases}$$
\end{przyklad}
\begin{przyklad}[Sfera w metryce dyskretnej]
$$S(x_0, r) = \{x\in X: \rho(x,x_0) = r\} = \begin{cases}
	\emptyset &\text{ gdy } r\in (0,1) \cup (1,\infty)\\
	X\smallsetminus\{x_0\} &\text{ gdy } r=1
\end{cases}$$
\end{przyklad}

\begin{df}[zbieżności ciągu o wyrazach w przestrzeni metrycznej]~\\
Ciąg $(a_n)$ o wyrazach w przestrzeni metrycznej $(X, \rho)$ jest zbieżny do $a\in X \Leftrightarrow \forall\varepsilon >0~~ \exists N ~ \forall n \geqslant N \quad \rho(a_n, a) < \varepsilon \Leftrightarrow \forall\varepsilon >0~~ \exists N ~ \forall n \geqslant N \quad |\rho(a_n, a) - 0| < \varepsilon \Leftrightarrow \rho(a_n, a) \rightarrow 0$
\end{df}

\begin{df}
Ciąg $(a_n)$ o wyrazach w przestrzeni metrycznej spełnia warunek Cauch'ego $\Leftrightarrow \forall \varepsilon >0 ~~ \exists N ~ \forall n,m > N \quad \rho(a_n, a_m) < \varepsilon$
\end{df}

\begin{tw}
Ciąg $(a_n)$ o wyrazach w przestrzeni metrycznej jest zbieżny $\Rightarrow$ spełnia warunek Cauch'ego
\end{tw}

\begin{przyklad}
Ciąg $a_n = \frac{1}{n}$ o wyrazach w przestrzeni metrycznej $X = (0, \infty)$ z metryką naturalną spełnia warunek Cauch'ego a mimo to nie jest zbieżny bo jedyny kandydat na granicę -- $0$ odpada
\end{przyklad}

\begin{df}[Przestrzeni metrycznej zupełnej]
Przestrzeń metryczna, w której każdy ciąg spełniający warunek Cauch'ego jest zbieżny nazywamy przestrzenią metryczną zupełną.
\end{df}

\begin{tw}
Przestrzeń $X = \mathbb{R}$ i $\rho (x,y) = |x-y|$ jest przestrzenią zupełną
\end{tw}

\begin{tw}[Banach o punkcie stałym]~\\
Niech $(X, \rho)$ będzie przestrzenią metryczną zupełną i $f: X \to X$ będzie odwzorowaniem zwężającym, to znaczy funkcją spełniająca warunek
$$ \exists L\in (0,1) ~~ \forall x,y\in X \quad \rho (f(x), f(y)) \leqslant L\rho (x, y) $$
Wówczas istnieje dokładnie jedno $x_0 \in X$ (nazywane punktem stałym odwzorowania), takie że $f(x_0) = x_0$. 
\end{tw}


\subsection{Elementy topologii}
\begin{df}[Zbioru otwartego]~\\
Zbiór $A$ zawarty w przestrzeni metrycznej $(X, \rho)$ nazywamy otwartym jeśli każdy punkt $a$ ze zbioru $A$ należy do $A$ wraz z pewną kulą o środku w $a$.\\
$$A \text{ jest otwarty } \Leftrightarrow \forall a\in A ~~ \exists r>0 \quad K(a, r) \subset A$$
\end{df}
\begin{df}[Zbioru domkniętego]~\\
Zbiór $A$ zawarty w przestrzeni metrycznej $(X, \rho)$ nazywamy domkniętym $ \Leftrightarrow X \smallsetminus A$ jest otwarty
\end{df}

\begin{df}
Wnętrze zbioru $A$ w przestrzeni metrycznej $(X, \rho)$ to 
$$ IntA := \{a\in A: ~~ \exists r>0 ~~ K(a,r) \subset A \} $$
\end{df}
\begin{df}
Domknięcie zbioru $A$ w przestrzeni metrycznej $(X, \rho)$ to 
$$ \overline{A} := \{x\in X: ~~ \exists r>0 ~~ K(x,r) \cap A \neq \emptyset \} $$
\end{df}

\begin{przyklad}
Rozpatrzmy przestrzeń metryczną $(X, \rho)$, gdzie $X = \mathbb{R}$, $\rho (x,y) = |x-y|$\\
$$Int([a,b)) = (a,b) \quad\quad \overline{(a,b)} = [a,b]$$
\end{przyklad}



\section{Funkcje wielu zmiennych}
\subsection{Granica funkcji}
\begin{df}[Granicy wg Heinego]
Mówimy, że $g$ jest granicą funkcji $f$ w punkcie $a$
$$\Leftrightarrow ~~\forall \{\mathbf{x}_n \subset D \smallsetminus \{a\} \} \quad [\lim\limits_{n\to\infty}\mathbf{x}_n = a \Rightarrow \lim\limits_{n\to\infty}f(\mathbf{x}_n)=g]$$
\end{df}
\begin{df}[Granicy wg Cauchego]
Mówimy, że $g$ jest granicą funkcji $f$ w punkcie $a$
$$\Leftrightarrow \forall\varepsilon >0 ~~\exists \delta >0 ~~\forall x\in D \quad 0<||\mathbf{x} - a|| <\delta \Rightarrow |f(\mathbf{x})-g| < \varepsilon$$
\end{df}
\begin{df}[Granicy niewłaściwej wg Heinego]
$$\lim\limits_{x\to a} = \pm \infty~~\Leftrightarrow ~~\forall \{\mathbf{x}_n \subset D \smallsetminus \{a\} \} \quad [\lim\limits_{n\to\infty}\mathbf{x}_n = a \Rightarrow \lim\limits_{n\to\infty}f(\mathbf{x}_n)=\pm\infty]$$
\end{df}
\begin{df}[Granicy niewłaściwej wg Cauchego]
$$\lim\limits_{x\to a} = +\infty~~\Leftrightarrow \forall M\in\mathbb{R}~~\exists \delta >0 ~~\forall x\in D \quad 0<||\mathbf{x} - a|| <\delta \Rightarrow f(x) > M$$
$$\lim\limits_{x\to a} = -\infty~~\Leftrightarrow \forall M\in\mathbb{R}~~\exists \delta >0 ~~\forall x\in D \quad 0<||\mathbf{x} - a|| <\delta \Rightarrow f(x) < M$$
\end{df}

\subsection{Ciągłość funkcji}
\begin{df}
Funkcja $f$ jest ciągła w punkcie $a\in D \Leftrightarrow [\exists \lim\limits_{x\to a}f(x) \text{ i } \lim\limits_{x\to a}f(x) = f(a)]$ 
\end{df}
\begin{df}
Funkcja $f$ jest ciągła $\Leftrightarrow$ jest ciągła w każdym punkcie $a\in D$
\end{df}

\begin{tw}[O ciągłości działań arytmetycznych]~\\
$[f,g: D\to \mathbb{R}, D\subset\mathbb{R}^k, a\in D; f,g \text{ są ciągłe w} a, x\in\mathbb{R}] \Rightarrow$ ciągłe w punkcie są następujące funkcje

\begin{itemize}
	\item $|f| $
	\item $\lambda f $
	\item $f \pm g $
	\item $fg$ 
	\item $\frac{f}{g} \text{ jeśli tylko } g(a) \neq 0$
\end{itemize}

\end{tw}
\begin{przyklad}
$ f:\mathbb{R}^2 \to\mathbb{R}, f(xy) \begin{cases} 
	\frac{xy}{x^2+y^2} & \text{ dla } (x,y) \neq (0,0)\\
	0 & \text{ dla } (x,y) = (0,0)
\end{cases}$\\

\begin{itemize}
	\item $f$ jest ciągła na $\mathbb{R}^2 \smallsetminus \{(0,0)\}$ jako iloraz funkcji ciągłych (wielomianów)
	\item $f$ nie jest ciągła w punkcie $(0,0)$, bo nie istnieje granica $\lim\limits_{(x,y)\to(0,0)}f(x,y)$
\end{itemize}
\end{przyklad}

\subsection{Pochodne i różniczkowalność funkcji wielu zmiennych}
\begin{df}[Pochodnej cząstkowej]~\\
Jeśli istnieje skończona granica $\lim\limits_{h\to 0} \frac{f(x_{10}+h, x_{20}, \dots, x_{k0})-f(x_{10}, x_{20}, \dots, x_{k0})}{h}$ to nazywamy ją pochodna cząstkową funkcji $f$ względem zmiennej $x_1$ w punkcie $x_0 = (x_{10}, x_{20}, \dots, x_{k0})$ i oznaczamy $f'_{x_1}(x_0)$		
\end{df}
\begin{df}[Pochodnej kierunkowej]~\\
Pochodną kierunkowa funkcji $f$ w punkcie $x_0 = (x_{10}, x_{20}, \dots, x_{k0})$ w kierunku wektora  $v = (v_{1}, v_{2}, \dots, v_{k})$ o długości 1 (to znaczy $||v|| = 1$) nazywamy granicę $\lim\limits_{t\to  0}\frac{f(x_0+tv) - f(x_0)}{t}$ jeśli istnieje i jest skończona. Oznaczamy ją $f'_v(x_0)$
\end{df}
\begin{df}[Różniczkowalności funkcji]~\\
Niech $G\subset \mathbb{R}^k$ będzie obszarem i $f:G\to \mathbb{R}$. Mówimy że $f$ jest różniczkowalna w punkcie $$x_0\in G \Leftrightarrow \exists A=(a_1, a_2, \dots, a_k) \in \mathbb{R}^k \quad \lim\limits
_{h\to 0}\frac{f(x_0+h)-f(x_0)-Ah}{||h||} = 0$$
\end{df}

%dodać definicję gradientu

\begin{tw}[Warunek konieczny różniczkowalności]~\\
$f:G\to\mathbb{R}, ~ G\subset\mathbb{R}^m$ to obszar, $f$ jest różniczkowalna w punkcie $x_0 \Rightarrow f$ jest ciągła w punkcie $x_0$
\begin{proof}
$f$ różniczkowalna w punkcie $x_0 \Rightarrow \lim\limits_{h\to 0}\frac{f(x_0+h)-f(x_0)-Ah}{||h||}=0$ dla pewnego $A=(a_1,a_2, \dots, a_m) \in \mathbb{R}^m$\\
Oznaczmy $\eta_{x_0}(h) = \frac{f(x_0+h)-f(x_0)-Ah}{||h||}$, wówczas mamy $\lim\limits_{n\to 0}\eta_{x_0}(h) = 0$\\
$f(x_0+h)=\eta_{x_0}(h) = \eta_x(h)||h|| + f(x_0) + Ah$\\
$\lim\limits_{h\to 0}f(x_0+h) = \lim\limits_{h\to 0}[\eta_{x_0}(h)] = \lim\limits_{h\to 0}[\underbrace{\eta_x(h)||h||}_{0} + f(x_0) + \underbrace{a_1h_1 + a_2h_2 + \dots + a_mh_m}_{0}] = f(x_0) \Rightarrow f$ jest ciągła w punkcie $x_0$
\end{proof}
\end{tw}

\begin{tw}[warunek dostateczny różniczkowalności]~\\
$f: G\to\mathbb{R}, ~ G\subset \mathbb{R}^m$ to obszar\\
Istnieją pochodne cząstkowe $f'_{x_1}(x)\dots f'_{x_m}(x)$ w pewnym otoczeniu punktu $x_0$ i są one ciągłe w punkcie $x_0 \Rightarrow f$ jest różniczkowalna w punkcie $x_0$
\end{tw}

\subsection{Pochodne cząstkowe wyższych rzędów}
\begin{df}
Niech $f: G\to\mathbb{R}$, gdzie $G\subset \mathbb{R}^m$ to obszar, ma pochodne cząstkowe $f'_{x_i}, i=1,\dots , m$\\
Jeśli funkcja $f'_{x_i}$ ma pochodną cząstkową po zmiennej $x_j$ to nazywamy ją pochodną czastkową drugiego rzędu i oznaczamy $f''_{x_ix_j} = (f'_{x_i})'_{x_j}$
\end{df}

\begin{tw}[Schwarza]~\\
Niech $f: G\to\mathbb{R}$, gdzie $G\subset \mathbb{R}^2$ to obszar.\\
Pochodne mieszane $f''_{xy}$ i $f''_{yx}$ istnieją w pewnym otoczeniu punktu $(x_0, y_0)\in G$ i są ciągłe w tym punkcie $\Rightarrow f''_{xy} = f''_{yx}$
\end{tw}

\begin{tw}[Wzór Taylora dla funkcje 2 zmiennych]~\\
Jeśli $f$ jest funkcją klasy $C^n$ w pewnym obszarze zawierającym odcinek $\overline{x_0x}$, to wewnątrz  tego odcinka znajduje się punkt $\mathbf{c} = (c_1,c_2)$ taki że\\
$
f(x) = f(x_0) + \frac{1}{1!}\left[\frac{df}{dx}(x_0)(x-x_0) + \frac{df}{dy}(x_0)(y-y_0)\right] +
				\frac{1}{2!}\left[\frac{d^2f}{dx^2}(x_0)(x-x_0)^2 + 2\frac{d^2f}{dxdy}(x_0)(x-x_0)(y-y_0) + \frac{d^2f}{dy^2}(x_0)(y-y_0)^2\right]
				+ \dots + \frac{1}{(n-1)!}\sum\limits_{i=0}^{n-1}{n \choose i}\frac{d^nf}{dx^idy^j}(x_0)(x-x_0)^i(y-y_0)^{n-i-1} + R_n(x_0) \text{, gdzie } R_n(x_0) = \frac{1}{n!}\sum\limits_{i=0}^n {n \choose i}\frac{d^n}{dx^idY^{n-1}}(c)(x-x_0)^i(y-y_0)^{n-1}
$
\end{tw}

\subsection{Ekstrema funkcji wielu zmiennych}
\begin{df}
Mówimy, że funkcja $f$ ma w punkcie $x_0\in G$:\\
\begin{itemize}
	\item maksimum lokalne jeśli $\exists r>0 ~\forall x\in K(x_0,r) \quad f(x) \leqslant f(x_0)$
	\item maksimum lokalne właściwe jeśli $\exists r>0 ~\forall x\in K(x_0,r), x \neq x_0  \quad f(x) < f(x_0)$
	\item minimum lokalne jeśli $\exists r>0 ~\forall x\in K(x_0,r) \quad f(x) \geqslant f(x_0)$
	\item minimum lokalne właściwe jeśli $\exists r>0 ~\forall x\in K(x_0,r), x \neq x_0  \quad f(x) > f(x_0)$
\end{itemize}
\end{df}

\begin{tw}[Warunek konieczny istnienia ekstremum funkcji wielu zmiennych]~\\
Jeśli istnieją pochodne cząstkowe $f'_{x_1}(x_0), f'_{x_2}(x_0), \dots, f'_{x_k}(x_0)$ i funkcja ma w punkcie $x_0$ ekstremum lokalne to $f'_{x_1}(x_0) = f'_{x_2}(x_0) = \dots = f'_{x_k}(x_0) = 0$
\end{tw}

\begin{tw}[Warunek dostateczny istnienia ekstremum dla funkcji 2 zmiennych]~\\
Niech $G\subset \mathbb{R}^2$ będzie obszarem, $f\in C^2(G)$ i $(x_0,y_0) \in G$. Jeśli $f'_x(x_0, y_0) = f'_y(x_0,y_0) = 0$ i $W(x_0,y_0) = 
\begin{vmatrix}
	f''_{xx}(x_0,y_0) & f''_{xy}(x_0,y_0)\\
	f''_{yx}(x_0,y_0) & f''_{yy}(x_0,y_0)\\
\end{vmatrix} > 0$ to w punkcie $(x_0,y_0)$ jest ekstremum lokalne właściwe. Ponadto jeśli $f''_{xx}(x_0,y_0)>0$ to jest to minimum lokalne, a jeśli $f''_{xx}(x_0,y_0)<0$ to maksimum lokalne
\end{tw}

\begin{df}
Formę kwadratową, a także odpowiadającą jej macierz $A$ nazywamy
\begin{itemize}
	\item dodatnio określoną $\Leftrightarrow \forall \mathbf{x}\neq 0 \quad \mathbf{x}A\mathbf{x}^T > 0$
	\item ujemnie określoną $\Leftrightarrow \forall \mathbf{x}\neq 0 \quad \mathbf{x}A\mathbf{x}^T < 0$
	\item nieujemnie określoną $\Leftrightarrow \forall \mathbf{x}\neq 0 \quad \mathbf{x}A\mathbf{x}^T \geqslant 0$
	\item niedodatnio określoną $\Leftrightarrow \forall \mathbf{x}\neq 0 \quad \mathbf{x}A\mathbf{x}^T \leqslant 0$	
	\item nieokreśloną $\Leftrightarrow \forall \mathbf{x,y} \quad \mathbf{x}A\mathbf{x}^T < 0,~ \mathbf{y}A\mathbf{y}^T > 0$
\end{itemize}
\end{df}

\begin{tw}[Kryterium Sylwestera]~\\
Macierz symetryczna $k\times k$:
\begin{itemize}
	\item jest dodatnio określona $\Leftrightarrow \forall i=1,\dots ,k ~ \det A^{(i)}>0$
	\item jest ujemnie określona $\Leftrightarrow \forall i=1,\dots ,k ~ (-1)^i\det A^{(i)}>0$
\end{itemize}
gdzie $A^{(i)}$ to macierz powstała z $A$ przez skreślenie kolumn i wierszy o numerach większych niż $i$
\end{tw}

\begin{tw}[warunek wystarczający istnienia ekstremum dla funkcji wielu zmiennych]~\\
Niech $G\subset \mathbb{R}^k$ będzie obszarem i niech $f\in C^2(G)$. Jeśli $\mathbf{x_0}\in G, f'x_0(\mathbf{x_0}) = f'x_1(\mathbf{x_0}) = \dots = f'x_n(\mathbf{x_0})$ to
\begin{itemize}
	\item Jeśli macierz $H(\mathbf{x_0})$ jest dodatnio określona, to $f$ ma minimum lokalne w $\mathbf{x_0}$
	\item Jeśli macierz $H(\mathbf{x_0})$ jest ujemnie określona, to $f$ ma maksimum lokalne w $\mathbf{x_0}$
	\item Jeśli macierz $H(\mathbf{x_0})$ jest nieokreślona, to $f$ nie posiada ekstremum lokalnego w $\mathbf{x_0}$
\end{itemize}
\end{tw}

\begin{tw}[O funkcji uwikłanej]~\\
Załóżmy że $F$ ma ciągłe pochodne cząstkowe w pewnym otoczeniu punktu $(x_0, y_0)$ oraz $F(x_0,y_0) = 0,~ F'_y(x_0,y_0) \neq 0$ to istnieje otoczenie $U_{x_0} \ni x_0$ i otoczenie $V_{y_0} \ni y_0$ oraz jedyna funkcja $y: U_{x_0} \to V_{y_0}$ klasy $C^1$ taka że $\forall x\in U_{x_0} \quad F(x, y(x)) = 0$. Ponadto $y$ ma ciągłą pochodną w $U_{x_0}$ i $y'(x) = \frac{-F_x(x,y)}{F'_y(x,y)}$, dla $x\in U_{x_0}$
\end{tw}



\section{Funkcje wektorowe}
\begin{df}
Niech $D\subset\mathbb{R}^m, ~ f_i: D\to\mathbb{R}, i=1,\dots ,k$ wtedy funkcją wektorową $m$ zmiennych nazywamy odwzorowanie $$f:D\to\mathbb{R}^k, f(x_1,\dots ,x_m) = (f_1(x_1,\dots ,x_m), \dots ,f_k(x_1,\dots ,x_m))$$
\end{df}

\begin{df}
Niech $G\subset\mathbb{R}^m$, będzie obszarem a $f: G\to\mathbb{R}^k$ funkcją wektorową $m$ zmiennych. Wówczas: 
\begin{itemize}
	\item $f$ ma granicę w punkcie $\mathbf{x_0} = (x_{10}, \dots ,x_{m0})$ który jest punktem skupienia zbioru $G$, równą $\mathbf{g} = (g_1, \dots , g_k) \Leftrightarrow \forall \varepsilon > 0 ~~ \exists \delta > 0 \forall )<||\mathbf{x}-\mathbf{x_0}||<\delta \Rightarrow ||f(\mathbf{x})-f(\mathbf{x_0})||<\varepsilon \Leftrightarrow \forall\{\mathbf{x_n}\}\subset G [\lim\limits_{n\to\infty}\mathbf{x_n} = \mathbf{x_0} \Rightarrow \lim\limits_{n\to\infty}f(\mathbf{x_n}) = f(\mathbf{x_0})$
	\item $f$ jest ciągła w punkcie $\mathbf{x_0}\in G \Rightarrow$ istnieje $\lim\limits_{x\to x_0}f(\mathbf{x}) = f(\mathbf{x_0})]$
\end{itemize}
\end{df}

\begin{df}[Różniczkowalności w punkcie]~\\
Niech $G\subset\mathbb{R}^m$, będzie obszarem a $f: G\to\mathbb{R}^k$ funkcją wektorową $m$ zmiennych. Mówimy, że $f$ jest różniczkowalna w punkcie $\mathbf{x}_0\in G \Leftrightarrow$ istnieje operator liniowy $A: \mathbb{R}^m \to \mathbb{R}^k$ taki że $\lim\limits_{h\to 0}\frac{f(\mathbf{x}_0+h)-f(\mathbf{x}_0)-Ah}{||h||} = 0$
\end{df}

\begin{df}
Krzywą w $\mathbb{R}^k$ nazywamy ciągłą funkcje wektorową jednej zmiennej $\gamma: J\to\mathbb{R}^k$. Jeśli funkcja jest różnowartościowa to wtedy krzywą nazywamy łukiem (zwykłym).
\end{df}
\begin{df}
Krzywą $\gamma: J\to\mathbb{R}^k$. Nazywamy łukiem gładkim, jeśli jest łukiem zwykłym i $\gamma_i, \dots ,\gamma_k$ są klasy $C^1$ oraz $\forall t\in J ||\gamma '(x)||^2 \neq 0$
\end{df}

\begin{tw}[Wzór na długość łuku gładkiego]~\\
Jeżeli krzywa $\gamma :[\alpha , \beta ] \to\mathbb{R}^2$ jest łukiem gładkim, to jest prostowalna i jej długość wyraża się wzorem $d = \int\limits_\alpha^\beta \sqrt{(x'(t))^2 + (y'(t))^2}dt$
\end{tw}

\section{Miara Jordana}
\begin{df}
Miara wewnętrzna Jordana zbioru $D$ to $\sup\{s:$ s to suma miar skończonej liczby kostek n-wymiarowych, takich że wnętrza tych kostek są parami rozłączne i każda z tych kostek zawiera się w zbiorze $D\} = \underline{j}(D)$
\end{df}
\begin{df}
Miara zewnętrzna Jordana zbioru $D$ to $\inf\{s:$ s to suma miar skończonej liczby kostek n-wymiarowych, pokrywających zbiór $ D\} = \overline{j}(D)$
\end{df}

\begin{df}
Zbiór $D$ jest mierzalny w sensie Jordana $\Leftrightarrow$ miara wewnętrzna Jordana zbioru $D = $ miara zewnętrzna Jordana zbioru $D$. 
Jeżeli $D$ jest mierzalny w sensie Jordana, to jego miara Jordana $j(D) = \underline{j}(D) = \overline{j}(D)$
\end{df}

\begin{df}
Obszar ograniczony $D\subset \mathbb{R}^2$ nazywamy obszarem regularnym jeśli jego brzeg da się rozbić na skończoną liczbę krzywych, które można zapisać w postaci $y=y(x), x\in [a,b]$ i $y\in C([a,b])$ lub $x=x(y), y\in [c,d]$ i $x\in C([c,d])$. Niektóre z tych krzywych mogą się redukować do punktu.
\end{df}
\begin{tw}
Każdy obszar regularny w $\mathbb{R}^2$ jest mierzalny w sensie Jordana.
\begin{proof}
Brzeg każdego obszaru regularnego w $\mathbb{R}^2$ ma miarę Jordana równą 0. Wynika to z definicji i tego że $y=f(x), x\in [a,b]$ i $y\in C([a,b])$ ma miarę Jordana w $\mathbb{R}^2$ równą $0$
\end{proof}
\end{tw}

\section{Całka Riemanna w $\mathbb{R}^n$}
\begin{df}
Ciąg podziałów $(\pi(n))$ nazywamy normalnym, jeśli $\delta (\pi (n)) \xrightarrow{n\to\infty}0$
\end{df}
\begin{df}
Jeśli istnieje stała $\delta\in\mathbb{R}$ taka, że dla dowolnego podziału $\pi_{n}$ normalnego prostokąta $P$ i dla dowolnego wartościowania $w_{n}$ mamy $\delta=\lim\limits_{n\to\infty}\delta_{n}(\pi_{n},w_{n}) $, to mówimy, że $f$ jest całkowalna w sensie Riemanna na $P$ i zapisujemy $\delta=\iint f(x,y)dxdy$, oraz $\delta$ nazywamy całką podwójną  funkcji $f$ na prostokącie $P$.
\end{df}
\begin{tw}
Dla każdego ciągu podziałów $\pi_{n}$ istnieją granice $\lim\limits_{n\to\infty}s_{n}(\pi_{n})$ i $\lim\limits_{n\to\infty}S_{n}(\pi_{n})$ i granice te zależą od wyboru ciągu podziałów: 
\begin{itemize}
	\item$s:=\lim\limits_{n\to\infty}s_{n}(\pi_{n})$ (całka podwójna górna),
	\item$S:=\lim\limits_{n\to\infty}S_{n}(\pi_{n})$(całka podwójna dolna)
\end{itemize}
\end{tw}

\end{document} 

